\documentclass[11pt,a4paper]{article}
\usepackage{amsmath,amssymb,amsthm}
\newtheorem{theorem}{Theorem}
\newtheorem{lemma}[theorem]{Lemma}
\begin{document}
\title{Irrationality of $\sqrt2$ and Beyond}
\author{Tomasz Wi\'sniewski}
\date{\today}
\maketitle
\section{The classical argument}
\begin{theorem}\label{thm:sqrt2}
The number $\sqrt2$ is irrational.
\end{theorem}
\begin{proof}
Suppose $\sqrt2=a/b$ with $\gcd(a,b)=1$. Then $a^2=2b^2$, so $a$ is even, say $a=2c$. Hence $2c^2=b^2$ and $b$ is even too, contradicting $\gcd(a,b)=1$.
\end{proof}

\section{Generalisation}
\begin{lemma}\label{lem:root}
If $n$ is a positive integer that is not a perfect square, then $\sqrt n$ is irrational.
\end{lemma}
\begin{proof}
Suppose $\sqrt n=a/b$ in lowest terms. Then $a^2=nb^2$. Any prime $p$ dividing $b$ would divide $a^2$ and hence $a$; so $b=1$ and $n=a^2$ is a perfect square.
\end{proof}

Lemma~\ref{lem:root} subsumes Theorem~\ref{thm:sqrt2}. Using the same idea one checks that
\begin{equation}
\sqrt[3]{2}\notin\mathbb{Q}.
\end{equation}
Hardy and Wright~\cite{hw} give many further proofs, and Niven~\cite[p.~21]{niven} treats the cubic case.

\begin{thebibliography}{9}
\bibitem{hw} G.~H. Hardy and E.~M. Wright, \emph{An Introduction to the Theory of Numbers}, Oxford University Press.
\bibitem{niven} I. Niven, \emph{Irrational Numbers}, Mathematical Association of America.
\end{thebibliography}
\end{document}
